%-------------------------
% Resume in Latex
% Author : Jake Gutierrez
% Based off of: https://github.com/sb2nov/resume
% License : MIT
%------------------------

\documentclass[letterpaper,11pt]{article}

\usepackage{latexsym}
\usepackage[empty]{fullpage}
\usepackage{titlesec}
\usepackage{marvosym}
\usepackage[usenames,dvipsnames]{color}
\usepackage{verbatim}
\usepackage{enumitem}
\usepackage[hidelinks]{hyperref}
\usepackage{fancyhdr}
\usepackage[english]{babel}
\usepackage{tabularx}
\usepackage{fontawesome5}
\input{glyphtounicode}


%----------FONT OPTIONS----------
% sans-serif
% \usepackage[sfdefault]{FiraSans}
% \usepackage[sfdefault]{roboto}
% \usepackage[sfdefault]{noto-sans}
% \usepackage[default]{sourcesanspro}

% serif
% \usepackage{CormorantGaramond}
% \usepackage{charter}


\pagestyle{fancy}
\fancyhf{} % clear all header and footer fields
\fancyfoot{}
\renewcommand{\headrulewidth}{0pt}
\renewcommand{\footrulewidth}{0pt}

% Adjust margins
\addtolength{\oddsidemargin}{-0.5in}
\addtolength{\evensidemargin}{-0.5in}
\addtolength{\textwidth}{1in}
\addtolength{\topmargin}{-.5in}
\addtolength{\textheight}{1.0in}

\urlstyle{same}

\raggedbottom
\raggedright
\setlength{\tabcolsep}{0in}

% Sections formatting
\titleformat{\section}{
  \vspace{-4pt}\scshape\raggedright\large
}{}{0em}{}[\color{black}\titlerule \vspace{-5pt}]

% Ensure that generate pdf is machine readable/ATS parsable
\pdfgentounicode=1

%-------------------------
% Custom commands
\newcommand{\resumeItem}[1]{
  \item\small{
    {#1 \vspace{-2pt}}
  }
}

\newcommand{\resumeSubheading}[4]{
  \vspace{-2pt}\item
    \begin{tabular*}{0.97\textwidth}[t]{l@{\extracolsep{\fill}}r}
      \textbf{#1} & #2 \\
      \textit{\small#3} & \textit{\small #4} \\
    \end{tabular*}\vspace{-7pt}
}

\newcommand{\resumeSubSubheading}[2]{
    \item
    \begin{tabular*}{0.97\textwidth}{l@{\extracolsep{\fill}}r}
      \textit{\small#1} & \textit{\small #2} \\
    \end{tabular*}\vspace{-7pt}
}


\newcommand{\resumeProjectHeading}[2]{
    \item
    \begin{tabular*}{0.97\textwidth}{l@{\extracolsep{\fill}}r}
      \small#1 & #2 \\
    \end{tabular*}\vspace{-7pt}
}

\newcommand{\resumeSubItem}[1]{\resumeItem{#1}\vspace{-4pt}}

\renewcommand\labelitemii{$\vcenter{\hbox{\tiny$\bullet$}}$}

\newcommand{\resumeSubHeadingListStart}{\begin{itemize}[leftmargin=0.15in, label={}]}
\newcommand{\resumeSubHeadingListEnd}{\end{itemize}}
\newcommand{\resumeItemListStart}{\begin{itemize}}
\newcommand{\resumeItemListEnd}{\end{itemize}\vspace{-5pt}}

%-------------------------------------------
%%%%%%  RESUME STARTS HERE  %%%%%%%%%%%%%%%%%%%%%%%%%%%%


\begin{document}

%----------HEADING----------
% \begin{tabular*}{\textwidth}{l@{\extracolsep{\fill}}r}
%   \textbf{\href{http://sourabhbajaj.com/}{\Large Sourabh Bajaj}} & Email : \href{mailto:sourabh@sourabhbajaj.com}{sourabh@sourabhbajaj.com}\\
%   \href{http://sourabhbajaj.com/}{http://www.sourabhbajaj.com} & Mobile : +1-123-456-7890 \\
% \end{tabular*}


\begin{tabular*}{\textwidth}{l@{\extracolsep{\fill}}r}
  \href{https://linkedin.com/in/proconlon/}{\textbf{\Large James Conlon}} &
  \href{https://github.com/proconlon}{\raisebox{-0.1\height}\faGithub\ github.com/proconlon} \\
  
  \small
  \href{mailto:proconlon52@gmail.com}{\raisebox{-0.1\height}\faEnvelope\ proconlon52@gmail.com}
  ~$\cdot$~
  \raisebox{-0.1\height}\faMapPin\ NYC Metro Area
  &
  \href{https://linkedin.com/in/proconlon/}{\raisebox{-0.1\height}\faLinkedin\ linkedin.com/in/proconlon} \\
\end{tabular*}

% \begin{center}
%     {\Huge \scshape James Conlon} \\    
%     \vspace{1pt}
%     \small
%     \raisebox{-0.1\height}\faLocationArrow\ NYC Metro Area~
%     \href{https://github.com/proconlon}{\raisebox{-0.2\height}\faGithub\ {github.com/proconlon}}~
%     \href{https://linkedin.com/in/proconlon/}{\raisebox{-0.2\height}\faLinkedin\ {linkedin.com/in/proconlon}}~
%     \href{mailto:proconlon52@gmail.com}{\raisebox{-0.2\height}\faEnvelope\  {proconlon52@gmail.com}}~ 
%     \vspace{-8pt}
% \end{center}


%-----------EDUCATION-----------
\section{Education}
  \resumeSubHeadingListStart
    \resumeSubheading
      {Boston University}{Boston, MA}
      {B.S. in Computer Engineering, Minor in Computer Science (GPA: 3.56)}{Sep 2022 -- May 2026}
  \resumeSubHeadingListEnd


%-----------EXPERIENCE-----------
\section{Experience}
  \resumeSubHeadingListStart

    \resumeSubheading
      {Embedded Software and Security Intern}{June 2025 -- Present}
      {ei$^3$ : Industrial Internet Intelligence}{Pearl River, NY}
      \resumeItemListStart
        \resumeItem{Designed and implemented a novel OpenVPN-based TAP server solution using Linux network namespaces, policy routing, and multiple ``virtual proxy'' VPN clients to resolve client-side NAT and subnet overlap issues for legacy industrial devices.}
        \resumeItem{Developed a comprehensive Python migration service to automate industrial router (Ewon, mGuard) configuration by downloading, validating, and uploading new profiles, ensuring zero-downtime migration while parsing and preserving full customer network topologies.}
        \resumeItem{Collaborated with internal support teams to author the SOP for automated remote deployments, establishing a safe migration process that preserves complex customer network topologies like VLANs and static routes.}
        \resumeItem{Designed and owned the REST API and data model (Python/Flask) for managing device configurations and firmware, allowing future expansion for new device models or new device-specific features.}
    \resumeItemListEnd
      
    \resumeSubheading
      {Embedded Software and Security Intern}{June 2024 -- June 2025}
      {ei$^3$ : Industrial Internet Intelligence \normalfont\small{\textit{(FT summer, PT during semesters)}}}{Pearl River, NY}
      \resumeItemListStart
        \resumeItem{Engineered a configurable OpenVPN TUN server in Docker, designed for deployment on edge hardware to establish a migration path for third-party industrial devices onto the ei$^3$ platform.}   
        % \resumeItem{Developed a multi-instance OpenVPN solution in Docker (Python, Bash), enabling advanced network routing (fwmark, policy routing, NAT) to secure legacy industrial devices in production.}
        % \resumeItem{Assisted with software rollout on hardware at customer locations by training support staff and documenting the onboarding processes for the product.}
        % \resumeItem{Integrated a containerized remote-access system with certificate management (OpenSSL) and IP allocation databases, ensuring future scalability.}
        % \resumeItem{Applied modern security best practices (Zero Trust, ISO/IEC frameworks) to safeguard older protocols and modernize industrial networks.}
        \resumeItem{Simulated factory networks (OPC-UA, MODBUS, PLCs), and performed competitive benchmarking against key competitors to evaluate technical capabilities, user experience, and security.}
      \resumeItemListEnd
      
% -----------Multiple Positions Heading-----------
%    \resumeSubSubheading
%     {Software Engineer I}{Oct 2014 - Sep 2016}
%     \resumeItemListStart
%        \resumeItem{Apache Beam}
%          {Apache Beam is a unified model for defining both batch and streaming data-parallel processing pipelines}
%     \resumeItemListEnd
%    \resumeSubHeadingListEnd
%-------------------------------------------

    \resumeSubheading
      {IT Support Specialist}{Aug 2022 -- Dec 2024}
      {BU Information Services and Technology}{Boston, MA}
      \resumeItemListStart
        \resumeItem{Diagnosed and resolved hardware/software issues: Windows/Mac/Linux, networking, and academic software.}
    \resumeItemListEnd

  \resumeSubHeadingListEnd


%-----------PROJECTS-----------
% \section{Projects}
%     \resumeSubHeadingListStart
%       \resumeProjectHeading
%           {\textbf{Gitlytics} $|$ \emph{Python, Flask, React, PostgreSQL, Docker}}{June 2020 -- Present}
%           \resumeItemListStart
%             \resumeItem{Developed a full-stack web application using with Flask serving a REST API with React as the frontend}
%             \resumeItem{Implemented GitHub OAuth to get data from user’s repositories}
%             \resumeItem{Visualized GitHub data to show collaboration}
%             \resumeItem{Used Celery and Redis for asynchronous tasks}
%           \resumeItemListEnd
%       \resumeProjectHeading
%           {\textbf{Simple Paintball} $|$ \emph{Spigot API, Java, Maven, TravisCI, Git}}{May 2018 -- May 2020}
%           \resumeItemListStart
%             \resumeItem{Developed a Minecraft server plugin to entertain kids during free time for a previous job}
%             \resumeItem{Published plugin to websites gaining 2K+ downloads and an average 4.5/5-star review}
%             \resumeItem{Implemented continuous delivery using TravisCI to build the plugin upon new a release}
%             \resumeItem{Collaborated with Minecraft server administrators to suggest features and get feedback about the plugin}
%           \resumeItemListEnd
%     \resumeSubHeadingListEnd

\section{Projects}
    \resumeSubHeadingListStart

    \resumeProjectHeading
      {\textbf{\underline{\href{https://github.com/proconlon/pm-edge-data-logger}{Predictive Maintenance Edge Data Logger}}} $|$ \emph{C, Python, AWS S3, Tailscale, BeagleBone}}{May 2025}
      \resumeItemListStart
        \resumeItem{Developed an end-to-end connected edge device data logger on a BeagleBone Black, using C for data acquisition from an OPC-UA server and Python for ML and cloud integration.}
        \resumeItem{Designed a dual-rate logging system to minimize cloud costs, capturing high-frequency data locally for ML model training and low-frequency data for long-term storage in AWS S3.}
        \resumeItem{Implemented secure remote access using Tailscale, air-gapping a factory network while enabling remote device management and data retrieval.}
        \resumeItem{Deployed a local Python ML model for predictive maintenance, analyzing high-frequency data to forecast equipment failures and automatically send email alerts to operators.}
      \resumeItemListEnd

      \resumeProjectHeading
        {\textbf{\underline{\href{https://github.com/proconlon/fpga-synth}{FPGA-Based Keyboard Synthesizer}}} $|$ \emph{Verilog, FPGA, Digital Signal Processing, Team Project}}{Dec 2023}
        \resumeItemListStart
          \resumeItem{Developed a digital synthesizer on Artix-7 FPGA, enabling simulation of musical notes and octaves.}
          \resumeItem{Designed a module for generating PWM signals with variable duty cycles to simulate sound waveforms including square, sine, and triangle waves.}
          \resumeItem{Integrated PS/2 keyboard input for control, along with a 7-segment display and audio output for feedback.}
          \resumeItem{Collaborated with team to address challenges such as sound generation and timing issues related to waveform indexing.}
        \resumeItemListEnd

     % \resumeProjectHeading
     %      {\textbf{\underline{\href{https://github.com/proconlon/RoomOccupancyMonitor}{Room Occupancy Monitor}}} $|$ \emph{Arduino, C++, Onshape, Team Project}}{Apr 2024}
     %      \resumeItemListStart
     %        \resumeItem{Collaborated on development of device to monitor the occupancy of a room to prevent overcrowding and create a mechanical barrier to entry.}
     %        \resumeItem{Primarily responsible for coding occupancy tracking, integrating various sensors (infrared, ultrasonic) and input/output devices to maximize accuracy and reliability.}
     %        \resumeItem{Developed a heuristic algorithm to manage ambiguous sensor data and varying traffic patterns, enhancing the accuracy of room occupancy monitoring.}
     %      \resumeItemListEnd

%    \resumeProjectHeading
%          {\textbf{Home Automation} $|$ \emph{Python, Networking, IoT, }}{Present}
%          \resumeItemListStart
%            \resumeItem{Developed and manage a home automation network, integrating over 10 IoT devices using MQTT and HTTP protocols for offline operation and security.}
%            \resumeItem{Proficient in selecting and configuring IoT hardware components including sensors (temperature, motion), smart switches, and lighting controls. Experience with microcontrollers (Arduino, Raspberry Pi) for device automation and control.}
%          \resumeItemListEnd
%          \vspace{-13pt}
            
    \resumeSubHeadingListEnd


%
%-----------PROGRAMMING SKILLS-----------
% \section{Technical Skills}
%  \begin{itemize}[leftmargin=0.15in, label={}]
%     \small{\item{
%      \textbf{Languages}{: Java, Python, C/C++, SQL (Postgres), JavaScript, HTML/CSS, R} \\
%      \textbf{Frameworks}{: React, Node.js, Flask, JUnit, WordPress, Material-UI, FastAPI} \\
%      \textbf{Developer Tools}{: Git, Docker, TravisCI, Google Cloud Platform, VS Code, Visual Studio, PyCharm, IntelliJ, Eclipse} \\
%      \textbf{Libraries}{: pandas, NumPy, Matplotlib}
%     }}
%  \end{itemize}

\section{Technical Skills}
 \begin{itemize}[leftmargin=0.15in, label={}]
    \small{\item{
     \textbf{Languages}{: Python, Go, C (systems-level), Bash, Verilog} \\
     \textbf{Framework \& Tools}{: Flask, Linux (Debian/Alpine/RHEL/Fedora), OpenSSL} \\
     \textbf{DevOps \& Virtualization}{: Docker, Docker Compose, Git, CI/CD (GitHub Actions), Jira} \\
     \textbf{Networking and Security}{: TCP/IP, OpenVPN, NAT, VLANs, Policy Routing, PKI, Wireshark} \\

     % \textbf{Courses}{: Distributed Systems, Embedded Systems, Operating Systems, Cybersecurity, Client-Server Software, Intro Databases} \\
     \textbf{Courses}{: {Systems: \{Distributed, Embedded, Operating\}}, Cybersecurity, Client-Server Software, Intro Databases} \\

    }}
 \end{itemize}
 \vspace{-16pt}

%-------------------------------------------
\end{document}

